\section{Conclusions and limitations}
\label{sec:conclusions}

A thin thermoelastic strip heated on its outer face bows \emph{away from the channel}: the through-wall
temperature difference sets a free curvature $\kth=(1+\nu)\alpha\qpp/k_s$ that is independent of the strip thickness,
and pinned ends turn it into a bow that widens the passage. At the baseline the reduced-order bow is
$\NumWscH\,H_0$ and a two-dimensional geometrically nonlinear finite-element solution of the wall gives
$\NLwFEH\,H_0$, which is $\NLoverScPct\,\%$ higher; the beam--column amplification of the free bow
($\NLampFE$ in the finite-element solution) and the relaxation of the membrane force
($\NLrelaxFE\,\%$ against $\NumRelax\,\%$) are both reproduced in magnitude. The binding limits on the heat flux are
deformation and stress, with buckling an outer ceiling. Indicative flow-side estimates, which need the coupled
computation for confirmation, point to a hydraulic effect of order one and a notable rise of the wall excess
temperature, while the effect of the flow back on the bow is small at the baseline.

The verification is partial. Only the wall has been verified (levels V1--V4 of Table~\ref{tab:hier}); the
fluid, the interface formulation and the coupled algorithm are specified and unverified (V5--V10).

\subsection*{Limitations}\label{sec:limits}
\begin{itemize}
\item \textbf{No coupled results.} Every flow-side number is a reduced-order estimate (Section~\ref{sec:indicative}),
and the lubrication and quasi-developed assumptions are not justified at the baseline parameters.
\item \textbf{Material data} are representative of an unfilled PEEK-class polymer and not frozen against a
datasheet; they are constant with temperature, and the strip is linearly thermoelastic (A3).
\item \textbf{Idealisations.} Perfect thermal contact (A6), ideal pinned ends (A7), an idealised uniform heater
(A8), and a casing that is only a support, with no contact with the deformed strip and no clearance (A9). A real
mounting would change the end conditions and, with them, the bow (Section~\ref{sec:q8free}).
\item \textbf{Heated face.} The model heats the outer face. Heating the fluid-facing face would bow the strip into the
channel and is a different model, outside this paper.
\item \textbf{Two dimensions.} Plane strain is assumed; end effects and finite thickness may explain part of the
finite-element departure from the beam model; shear flexibility accounts for a small share of it, and most of the
$\NLoverScPct\,\%$ has not been attributed.
\item \textbf{Steady state.} Transient heating, thermal inertia and viscoelasticity of the polymer are not
modelled.
\end{itemize}

\subsection*{Further work}
Carry out V5--V10 of Table~\ref{tab:hier}; compute the coupled problem and test the contraction of the
fixed point; replace the idealised heater and end conditions with a specified mounting and casing contact; and
validate against experiment.
